\subsection{Related Work}\label{sec:related-work}
\textbf{Prompting persona-like behaviour.}
Persona-like behaviour can be induced through interventions at different points in the model pipeline. The most lightweight way to shape a model persona is through natural-language prompts, which can produce large behavioural changes without modifying model weights~\citep{askell2021general, brown2020language, ouyang2022instructgpt, bai2022constitutional}. Persona prompting and role-play have also been studied as ways of inducing distinct behavioural profiles in LLMs~\citep{deshpande2023toxicity, park2023generative, shanahan2023role}. However, prompting is often brittle: the induced persona may weaken over long rollouts, conflict with later context, or drift as the conversation evolves.

\textbf{Modifying personas at inference time.}
A second class of interventions shapes personas by modifying the model's computation during inference, rather than by changing the prompt or updating the weights~\citep{turner2023activation, zou2023representation, rimsky2024steering, arditi2024refusal, chen2025persona}, and can be more targeted than prompting because they intervene on internal representations associated with particular traits. For persona control, activation capping has been proposed as a way to limit deviations from an assistant-like default~\citep{lu2026assistant}. Nevertheless, steering methods can depend sensitively on layer choice, prompt distribution, rollout length, and the stability of the relevant representation~\citep{tan2024analysing}.

\textbf{Training and post-training to shape model personas.}
A third class of interventions changes the model's weights. Pre-training exposes models to many human and fictional persona-types~\citep{tice2026alignmentpretraining}, while post-training biases this broad simulator toward narrower behavioural defaults, such as helpful, harmless, instruction-following assistant behaviour~\citep{ouyang2022instructgpt, bai2022constitutional, lu2026assistant}. More targeted character-training pipelines fine-tune or adapt models toward particular persona archetypes~\citep{maiya2025opencharacter, li2025big5}.

\textbf{Hybrid approaches which train model modifiers that can be applied at inference time.} A similar approach to ours is PsychAdapter~\citep{vu2026psychadapter}, which trains LoRAs together with per-layer projections that inject a tweakable vector of trait scores into every layer as a dummy token.

\textbf{Task arithmetic is a prior approach for model weight modification.} This method has shown that weight-space directions encode task-specific knowledge and can be composed additively~\citep{ilharco2023editing,huang2024chatvectorsimpleapproach}. We build on these methods by training trait-specific LoRAs and studying whether they provide a more flexible and robust way to construct personas than prompt-only or activation-only control.
